\scanpage{83}
\mspage{79}
% archive sheet number: 82

En fait, on a \inserted{donc} construit une application
\diagram{\begin{tikzcd}[ampersand replacement=\&]
\Sigma \arrow[r,phantom,"{\overset{\unsure{\beta|B}}{\unsure{\longrightarrow|{-}}}}"{description}] \& \Sigma_E
\end{tikzcd}}
visiblement croissante, telle que
\[
\unsure{\alpha|\varphi}\beta=\operatorname{id}_{\Sigma}.
\]
De plus, on vient de voir qu'elle est compatible \unsure{à|:} $\unsure{\complement|C|\ell|l}$. Pour achever la dém, il suffit de voir que
\[
\inserted{\beta\alpha}\deleted{\unsure{\beta\alpha\beta|\beta\alpha}}=\operatorname{id}_{\Sigma_E},
\]
car alors $\alpha$ et $\beta$ sont des \deleted{applications \unsure{inv.}}\inserted{iso. d'ens ordonnés} \deleted{\unsure{inv.} l'une de l'autre, \unsure{croi.}}\inserted{inverses l'une de l'autre, et (comme} on voit sur $\beta$) compatibles \unsure{à|:} $\complement$.

Soit donc $\unsure{A|\mathbb{A}|\mathscr{A}}\in\Sigma_E$, $A'=\unsure{\operatorname{comp}^{\circ}|\operatorname{comp}^{o}}(A\unsure{)}$, alors les $\sigma(x)$ et $\sigma(x')$ $(x\in A,\ x'\in A')$ sont disjoints\inserted{(par a))}, donc leurs $\operatorname{Sup}$ $\alpha(A)$, $\alpha(A')$ le sont aussi\deleted{\unsure{(par a))}}. \deleted{Considérons}
\[
\deleted{\beta\alpha(A)=\{x\in E\mid\sigma(x)\leq\alpha(A)\}}
\quad\deleted{\text{\unsure{Consid il s'ensuit}}}
\]
i.e.
\[
\alpha(A)\leq\complement\alpha(A'),\quad\text{d'où}
\]
\deleted{\unsure{on a}}
\[
\beta(\alpha(A))\leq\beta\complement\alpha(A')=\complement\beta(\alpha(A')),\quad\text{i.e.}
\]
\diagram{\begin{tikzcd}[ampersand replacement=\&,column sep=small,row sep=small]
\beta\alpha(A) \arrow[d,phantom,"\leq"{description,rotate=90}] \& \mathrel{|{\circ}|} \& \beta\alpha(A'). \arrow[d,phantom,"\leq"{description,rotate=90}] \\
A \& \& A'
\end{tikzcd}}
Mais cela implique
\[
A=\beta\alpha(A),\quad A'=\beta\alpha(A'),
\]
par le

\uline{Lemme} Soit $\Sigma$ un \unsure{spotos|spatos} ($\Sigma_E$ en \unsure{l'occurrence}), soient $A\leq B$, $A'\leq B'$ avec $\deleted{\unsure{A}}B\mathrel{|{\circ}|}B'$, et $A'=\complement A$. Alors $A=B$, $A'=B'$.

En effet, on a $A\leq B\leq\complement B'\leq\complement A'=A$, donc \unsure{tous les|toutes les} $\leq$ sont des \unsure{égalités}.
